\chapter{Metodologia}
\label{cap:metodologia}

\section{Classificação da pesquisa}
\label{sec:classificacao}

De acordo com a classificação de \citeonline[cap.~4]{gil2017}, esta pesquisa é aplicada quanto à finalidade, porque
parte de um problema concreto do aplicativo Bem-Te-Vi e gera artefatos destinados à sua resolução; além disso, 
trata-se de uma pesquisa qualitativa
quanto à abordagem, pois apresenta seus resultados por meio de descrições verbais, e não em termos numéricos.
As contagens de campos, estruturas e requisitos mencionadas no texto descrevem os artefatos e não são objeto
de análise estatística.

No que diz respeito ao procedimento, o trabalho segue a metodologia de \textit{design science} de
\citeonline[p.~3]{wieringa2014design}, que considera a pesquisa na área de engenharia de software como o projeto
e o estudo de artefatos inseridos no contexto do problema. Wieringa organiza
esse trabalho em um ciclo de projeto com três tarefas: examinar o problema, projetar um tratamento e validá-lo.
A validação, nessa metodologia,
ocorre antes da implementação e investiga se o artefato, uma vez inserido no contexto,
contribuiria para os objetivos das partes interessadas \cite[p.~27,~31]{wieringa2014design}.

Neste trabalho, o contexto é o acompanhamento das avaliações do projeto v-FES
cycling e os artefatos são a especificação de requisitos, a classificação dos níveis de automação
das escalas e o modelo de dados em HL7 FHIR. As Seções~\ref{sec:inventario} a~\ref{sec:automacao} correspondem à investigação do problema, e as
Seções~\ref{sec:especificacao} e~\ref{sec:modelo-fhir}, ao projeto. Como, nesta fase, não há reuniões com
a equipe de fisioterapia, a validação com as partes interessadas é adiada para a próxima fase do trabalho, juntamente com as decisões que exigem julgamento clínico. Neste trabalho, os artefatos
passam por verificação: o modelo é verificado quanto à cobertura dos formulários e testado com
instâncias de exemplo, de acordo com os critérios indicados na Seção~\ref{sec:criterios}.

\section{Fontes}
\label{sec:fontes}

Sem sessões de elicitação, os requisitos e o modelo partem das fontes documentais do Quadro~\ref{qua:fontes}. O
código do aplicativo é analisado no commit \texttt{a9059f0}, de 4 de novembro de 2026, da branch
\texttt{feat/cif-avaliacoes} \cite{higa2026bemtevi}. As afirmações sobre o aplicativo indicam o arquivo e a linha
desse commit, e cada requisito do Capítulo~\ref{cap:resultados} indica a fonte de onde foi derivado. Os caminhos de
arquivo são relativos ao repositório do aplicativo. Os arquivos em Python estão em \arq{Backend/} e
\arq{scripts/cif/}, e os arquivos em Dart, em \arq{lib/}. Os caminhos que começam em \arq{entregas/} são do
repositório do TCC.

\begin{quadro}[htb]
	\centering
	\caption{Fontes documentais do trabalho}
	\label{qua:fontes}
	\footnotesize
	\begin{tabularx}{\textwidth}{|>{\raggedright\arraybackslash}p{0.2\textwidth}|>{\raggedright\arraybackslash}X|>{\raggedright\arraybackslash}p{0.27\textwidth}|}
		\hline
		Fonte & Conteúdo & Uso (seção) \\ \hline
		Código do aplicativo & Telas em Flutter (\texttt{lib/}) e API em FastAPI (\texttt{Backend/})
			\cite{higa2026bemtevi} & Estado atual, regras implementadas e problemas de dados
			(\ref{sec:inventario}, \ref{sec:modulo-cif}) \\ \hline
		Esquema do banco & \texttt{Backend/schema.sql} e migração \texttt{001\_create\_avaliacoes\_cif.sql}
			& Ponto de partida do modelo FHIR (\ref{sec:inventario}, \ref{sec:modelo-fhir}) \\ \hline
		Planilha CIF da equipe & Abas Funções do Corpo, Estrutura do Corpo, Atividades e Participação, Fatores
			Ambientais, Plan1 e Legenda & Origem do catálogo e das regras de cálculo (\ref{sec:modulo-cif}) \\ \hline
		Documentos do módulo CIF & Catálogo (\texttt{catalogo.v0.1.json}), auditoria (\texttt{auditoria.v0.1.json},
			\texttt{regras-e-auditoria.md}) e plano do módulo (\texttt{docs/planos/cif.md}) & Regras verificáveis,
			achados nas fórmulas, decisões pendentes e etapas do módulo (\ref{sec:modulo-cif}) \\ \hline
		Protótipo & Arquivo no Figma e capturas das telas da CIF & Fluxo da avaliação CIF (\ref{sec:modulo-cif}) \\ \hline
		Plano de trabalho de iniciação científica & \citeonline{ochoadiaz2026plano} & Necessidades da base e dados
			para o módulo de IA (\ref{sec:especificacao}) \\ \hline
		Padrões de interoperabilidade & FHIR R4, SDC, BR-Core e guia da RNDS
			\cite{rnds2026guia} & Recursos, perfis e requisitos de interoperabilidade
			(\ref{sec:especificacao}, \ref{sec:modelo-fhir}) \\ \hline
		Normas & Portaria 1.434/2020, LGPD e Resolução COFFITO 414/2012
			(\citeauthor{lgpd2018}, \citeyear{lgpd2018}, \citeyear{portaria1434:2020}; \citeauthor{coffito414:2012}, \citeyear{coffito414:2012}) & Requisitos normativos, de privacidade e de
			registro (\ref{sec:especificacao}) \\ \hline
		Literatura dos instrumentos & Referências apresentadas no Capítulo~\ref{cap:referencial} & Regras de
			cálculo das escalas (\ref{sec:automacao}) \\ \hline
	\end{tabularx}
	\fonte{Elaborado pelo autor.}
\end{quadro}

Os exemplos e as instâncias de teste usam somente dados fictícios. Nenhum dado real de paciente é coletado ou
analisado nesta etapa.

\section{Etapas do trabalho}
\label{sec:etapas}

O trabalho segue sete etapas, agrupadas segundo as tarefas do ciclo de projeto na Figura~\ref{fig:fluxo}. Cada etapa
usa o resultado das anteriores e corresponde a uma seção do Capítulo~\ref{cap:resultados}. O levantamento do estado
atual e a análise do módulo CIF mostram o que o aplicativo já faz e onde estão os problemas de dados. A classificação
da automação indica o que cada escala permite calcular sem julgamento clínico. Esses três resultados alimentam a
especificação de requisitos, que orienta o modelo em HL7 FHIR. O modelo é então verificado com instâncias de exemplo,
e a consolidação reúne a rastreabilidade e as pendências que seguem para o TCC 2. Ao longo do
Capítulo~\ref{cap:resultados}, um item da CIF, b167 (``Funções mentais da linguagem''), serve de exemplo condutor:
ele é acompanhado da planilha da equipe até o modelo em HL7 FHIR e a matriz de rastreabilidade.

\begin{figure}[H]
	\centering
	\begin{tikzpicture}[
			font=\footnotesize,
			etapa/.style={draw, rounded corners=2pt, align=center, text width=3.6cm, minimum height=1.05cm},
			titulo/.style={font=\footnotesize\bfseries, align=center, text width=3.6cm},
			grupo/.style={draw, dashed, rounded corners=4pt, inner sep=5pt},
			seta/.style={-{Stealth[length=2mm]}, semithick},
			node distance=4mm and 12mm,
		]
		% Investigação do problema
		\node[titulo] (t1) {Investigação do problema};
		\node[etapa, below=2mm of t1] (e1) {1. Estado atual do aplicativo\\(Seção~\ref{sec:inventario})};
		\node[etapa, below=of e1] (e2) {2. Análise do módulo CIF\\(Seção~\ref{sec:modulo-cif})};
		\node[etapa, below=of e2] (e3) {3. Nível de automação das escalas\\(Seção~\ref{sec:automacao})};
		\path (e1) -- (e2) coordinate[midway] (m12);
		\path (e2) -- (e3) coordinate[midway] (m23);

		% Projeto do tratamento
		\node[etapa, anchor=west] (e4) at ([xshift=12mm]e1.east |- m12) {4. Especificação de requisitos\\(Seção~\ref{sec:especificacao})};
		\node[etapa, anchor=west] (e5) at ([xshift=12mm]e1.east |- m23) {5. Modelo de dados em HL7 FHIR\\(Seção~\ref{sec:modelo-fhir})};
		\node[titulo] (t2) at (e4 |- t1) {Projeto do tratamento};

		% Verificação
		\node[etapa, right=of e4] (e6) {6. Verificação com instâncias de exemplo\\(Seção~\ref{sec:modelo-fhir})};
		\node[etapa, right=of e5] (e7) {7. Consolidação: rastreabilidade e pendências\\(Seção~\ref{sec:consolidacao})};
		\node[titulo] (t3) at (e6 |- t1) {Verificação};

		% Molduras das tarefas, com a mesma altura
		\node[grupo, fit=(t1)(e1)(e2)(e3)] (g1) {};
		\node[grupo, fit=(t2)(e4)(e5)(e4 |- e3.south)] (g2) {};
		\node[grupo, fit=(t3)(e6)(e7)(e6 |- e3.south)] (g3) {};

		% Fontes
		\coordinate (fa) at ([yshift=8mm]g1.north west);
		\coordinate (fb) at ([yshift=16mm]g3.north east);
		\node[draw, rounded corners=2pt, fit=(fa)(fb), inner sep=0pt] (fontes) {};
		\node at (fontes) {Fontes documentais (Quadro~\ref{qua:fontes})};

		% Fluxo
		\draw[seta] (e1) -- (e2);
		\draw[seta] (e2) -- (e3);
		\draw[seta] (e4) -- (e5);
		\draw[seta] (e6) -- (e7);
		\draw[seta] (g1) -- (g2);
		\draw[seta] (g2) -- (g3);
		\foreach \g in {g1, g2, g3} {
			\draw[seta] (fontes.south -| \g.north) -- (\g.north);
		}
	\end{tikzpicture}
	\caption{Etapas do trabalho, agrupadas pelas tarefas do ciclo de projeto}
	\label{fig:fluxo}
	\fonte{Elaborada pelo autor com base em \citeonline{wieringa2014design}.}
\end{figure}

\section{Técnicas}
\label{sec:tecnicas}

\textbf{Engenharia reversa do código e do esquema.} \citeonline{chikofsky1990reverse} definem a engenharia reversa
como a análise de um sistema com o objetivo de identificar seus componentes e as relações entre eles, bem como de
representá-lo em um nível mais elevado de abstração. Na avaliação do estado atual (Seção~\ref{sec:inventario}), essa
análise percorre, em cada formulário, a tela em Flutter, as rotas da API e as tabelas do arquivo \texttt{schema.sql},
e registra a estrutura dos campos, as funções disponíveis e o comportamento que afeta a qualidade dos dados.

\textbf{Recuperação de requisitos e regras de negócio.} O código também serve como fonte de requisitos:
\citeonline{yu2005reverse} extraem modelos de objetivos a partir de código legado, e \citeonline{sneed1996extracting}
descrevem a extração de regras de negócio a partir do código-fonte. Como o módulo CIF foi implementado antes da
criação de uma especificação formal, ele é lido da mesma maneira (Seção~\ref{sec:modulo-cif}), e cada comportamento
implementado, como a média das notas dos itens filhos ou o estado pendente em caso de resposta ausente, é registrado
como um requisito ou regra de negócio, com indicação do arquivo que o implementa. Os cálculos dos demais formulários,
levantados na Seção~\ref{sec:inventario}, passam pelo mesmo procedimento. Como a obtenção de um requisito a partir do
código-fonte pode reproduzir erros da implementação, o requisito deve seguir o instrumento ou a norma sempre que o
comportamento implementado divergir da regra publicada para o instrumento ou de uma norma, e essa divergência é
registrada como um achado.

\textbf{Confronto da planilha com o catálogo.} O catálogo da CIF foi extraído da planilha da equipe, e a análise
compara, aba por aba, a hierarquia, os campos e as fórmulas da planilha com o catálogo versionado e com a auditoria
registrada no repositório do aplicativo. Cada divergência é classificada como decidida, caso já exista uma decisão
registrada, ou como pendente, caso dependa da confirmação da equipe.

\textbf{Derivação de requisitos a partir de normas e padrões.} \citeonline{breaux2008analyzing} propõem um método
para extrair de textos normativos direitos, obrigações e restrições, mantendo cada requisito vinculado ao trecho do
qual se origina. Este trabalho não aplica a análise completa proposta pelos autores, mas apenas a parte que garante
essa ligação: cada requisito derivado da LGPD, da Resolução COFFITO 414/2012, da Portaria 1.434/2020 ou do BR-Core
indica o artigo ou o perfil do qual provém.

\textbf{Classificação da automação.} Cada valor produzido pelos instrumentos, seja um item, um subtotal ou um escore,
recebe um dos quatro graus definidos na Seção~\ref{sec:criterios}: o grau atual baseia-se no código do aplicativo, e
o grau possível baseia-se na regra publicada para o instrumento em questão. Essa classificação é proposta neste
trabalho e baseia-se na distinção feita por \citeonline[p.~288]{parasuraman2000automation} entre aquisição de
informação, análise de informação e seleção de decisão. Os dois graus intermediários correspondem à automação da
análise, que envolve a integração de várias variáveis em um único valor, e o último grau indica o ponto em que a
decisão cabe ao profissional. É considerado valor cada item do instrumento e cada dado que entra no cálculo por ele
determinado, como o peso e a altura no IMC e a escolaridade na pontuação esperada do MEEM, enquanto identificação,
datas, nome do examinador, equipamento e texto livre não são levados em consideração. Na GAS, a contagem considera as
13 metas que o formulário sempre apresenta (\arq{gasForm.dart}, linhas 147--231), e, na CIF, são incluídos os itens
da cadeia de fórmulas de cada sexo (Seção~\ref{sec:cif-catalogo}).

\textbf{Conferência dos cálculos.} Cada valor calculado pelo aplicativo é comparado com a regra descrita pela fonte
do instrumento (Seção~\ref{sec:ref-instrumentos}). Os cálculos do servidor são executados sobre casos fictícios por
meio de scripts que importam o próprio código do commit analisado (\arq{conferir_calculo_gas.py} e
\arq{conferir_calculo_asia.py}, no diretório \arq{entregas/scripts/}), enquanto os cálculos realizados apenas na
tela, escritos em Dart, são transcritos para o script ou verificados durante a leitura do código. Na ASIA, o
resultado esperado de cada caso é obtido manualmente com base nas regras de
\citeonline[p.~3,~16,~18--19,~22]{rupp2021isncsci} e conferido com o algoritmo público do Praxis Spinal Cord
Institute, que implementa a revisão de 2019 do padrão \cite{praxis2023isncsci}. Se o código divergir da regra
publicada, essa divergência recebe um identificador DC, distinto dos problemas de dados (PD), pois o PD5 compara a
tela com o servidor, enquanto os DC comparam o aplicativo com o instrumento.

\textbf{Mapeamento para o HL7 FHIR.} Cada estrutura dos formulários é mapeada para recursos do FHIR R4
\cite{hl7fhirR4}, com preferência pelos perfis do BR-Core, quando existentes \cite{brcore2026}, e os cálculos são
declarados por meio de expressões do SDC \cite{hl7sdc}. A CIF é mapeada detalhadamente, e os demais formulários, de
acordo com sua estrutura.

\textbf{Verificação com instâncias de exemplo.} Para a CIF, são criadas instâncias em JSON com dados fictícios,
validadas em relação à especificação R4 e aos perfis do BR-Core utilizados no modelo. Neste trabalho, validar uma
instância significa conferir sua conformidade estrutural com a especificação e com os perfis, ou seja, os tipos, a
cardinalidade e a correspondência entre as respostas e o questionário. A correção dos valores calculados é verificada
separadamente: as expressões de cálculo do SDC são executadas sobre as instâncias, e os resultados são comparados com
os valores esperados, obtidos a partir da calculadora do aplicativo. As mesmas instâncias servem para conferir se
cada problema de dados da Seção~\ref{sec:inventario} tem uma resposta no modelo, enquanto os demais formulários,
conforme o recorte do trabalho, são analisados pela estrutura, sem instâncias.

\section{Critérios para responder às questões de pesquisa}
\label{sec:criterios}

\begin{description}
	\item [P1] Cobertura do modelo. Em cada formulário, as estruturas de dados (bloco repetido, campo calculado,
	variação por sexo, texto livre e valor ausente) são classificadas como representadas, se um recurso ou elemento
	padrão do FHIR for suficiente; representadas com adaptação, se exigirem extensão, perfil ou sistema de códigos
	próprio; ou não representadas. A resposta à P1 é afirmativa se nenhuma estrutura ficar sem representação e, na
	CIF, se as instâncias de exemplo contiverem todos os campos do formulário. As adaptações indicam onde o modelo
	depende de elementos próprios do projeto.
	\item [P2] Qualidade e rastreabilidade dos requisitos. Cada requisito funcional (RF), requisito não funcional
	(RNF) ou regra de negócio (RN) recebe identificador, descrição, fonte do Quadro~\ref{qua:fontes}, justificativa,
	prioridade e critério de aceitação. Os requisitos atendem às características de um requisito bem formulado
	segundo a ISO/IEC/IEEE 29148, como ser necessário, não ambíguo e verificável \cite{iso29148:2018}. O autor atribui
	a prioridade em três níveis: alta, para requisitos exigidos por norma, já implementados no módulo CIF ou
	necessários para corrigir problemas de dados da Seção~\ref{sec:inventario} ou divergências de cálculo da
	Seção~\ref{sec:automacao}; média, para requisitos sem os quais o acompanhamento entre avaliações não funciona ou
	perde confiabilidade, como histórico, estado, versão do formulário, comparação, regras de cálculo e de registro
	dos valores comparados e representação dos dados no padrão de interoperabilidade adotado; e baixa, para
	requisitos que ampliam o uso dos dados além do acompanhamento, como a exportação para fins de pesquisa. A equipe
	deve revisar essa priorização no TCC 2. A resposta à P2 estará completa quando cada requisito tiver fonte e
	critério de aceitação e estiver listado na matriz de rastreabilidade.
	\item [P3] Grau de automação. Cada valor das escalas é classificado em um dos quatro graus: digitado, se for um
	dado de entrada observado ou medido pelo profissional; calculado pelo aplicativo, quando o código já o obtém a
	partir de outros valores; calculável pela regra do instrumento, se a literatura definir uma regra determinística
	que o aplicativo ainda não implementa; e dependente de julgamento clínico, se o valor decorrer de outros dados,
	mas exigir a interpretação do profissional. O segundo grau é comprovado por arquivo e linha de código, e o
	terceiro, pela referência do instrumento. Os valores do terceiro grau representam as oportunidades de automação.
	Os do segundo grau são, além disso, conferidos com a regra publicada do instrumento, e cada divergência
	recebe um identificador DC (Seção~\ref{sec:tecnicas}). A resposta a P3 é a distribuição dos valores pelos
	quatro graus, com as divergências encontradas na conferência. Um terceiro grau vazio também é uma resposta
	possível e indica que, dentro de uma avaliação, o aplicativo já calcula todos os valores que os instrumentos
	definem por regra determinística.
\end{description}

\section{Ferramentas}
\label{sec:ferramentas}

\begin{itemize}
	\item Git, para versionar o texto e os artefatos. O código do aplicativo entra no repositório do trabalho como
	submódulo, fixado no commit analisado;
	\item Python, em scripts de conferência, como a comparação entre a planilha e o catálogo da CIF;
	\item Node.js, para executar o algoritmo do ISNCSCI do Praxis Spinal Cord Institute \cite{praxis2023isncsci}, com o
	qual os casos fictícios da ASIA são classificados e comparados com o servidor do aplicativo;
	\item validador oficial do HL7 FHIR, executado em Java, para validar as instâncias de exemplo;
	\item fhirpath.js \cite{hl7fhirpathjs}, executado em Node.js, para executar as expressões de cálculo do SDC sobre as
	instâncias e comparar os resultados com o calculador do aplicativo;
	\item Figma, para consultar o protótipo das telas.
\end{itemize}

Os scripts, as instâncias em JSON e as tabelas completas ficam no repositório do trabalho, na pasta
\texttt{entregas/}.

\section{Cronograma}
\label{sec:cronograma}

O Quadro~\ref{qua:cronograma} apresenta o cronograma do TCC 1, com entrega em 1º de dezembro de 2026. O módulo CIF
foi implementado antes, como parte do trabalho no projeto, e aparece na primeira linha por ser o principal resultado
parcial desta etapa. Uma versão completa do texto vai para a orientadora em 24 de novembro, e a última semana fica
reservada às correções.

\begin{quadro}[htb]
	\centering
	\caption{Cronograma do TCC 1}
	\label{qua:cronograma}
	\footnotesize
	\begin{tabularx}{\textwidth}{|p{3.2cm}|X|}
		\hline
		Período (2026) & Atividades \\ \hline
		até 30/09 & Implementação do módulo CIF no aplicativo: cálculo, persistência, formulário, histórico e
			comparação \\ \hline
		01/10 a 15/10 & Pré-TCC, levantamento da bibliografia e estudo dos formulários do aplicativo \\ \hline
		16/10 a 31/10 & Ajustes no módulo CIF, revisão dos demais módulos e correção de falhas \\ \hline
		01/11 a 07/11 & Recorte do tema, esquema do banco, auditoria da CIF e guia de implantação \\ \hline
		08/11 a 14/11 & Introdução, reunião de orientação de 12/11, que redefiniu o escopo, e estrutura do trabalho \\ \hline
		15/11 a 16/11 & Metodologia e Seção~\ref{sec:modulo-cif} \\ \hline
		17/11 a 20/11 & Seções~\ref{sec:inventario}, \ref{sec:automacao} e~\ref{sec:especificacao}, instâncias em HL7
			FHIR e parte do referencial teórico \\ \hline
		21/11 a 23/11 & Seções~\ref{sec:modelo-fhir} e~\ref{sec:consolidacao}, restante do referencial teórico,
			considerações finais e resumo \\ \hline
		24/11 & Envio da versão completa à orientadora \\ \hline
		25/11 a 01/12 & Correções indicadas pela orientadora, conferência das normas ABNT e entrega \\ \hline
	\end{tabularx}
	\fonte{Elaborado pelo autor.}
\end{quadro}
